% !TeX program = lualatex
% ================================================================
% Polish Tier 3 Vocabulary Version of FractionsQuickQuestions
%
% Generated by the server using the following settings:
%
%   language            Polish (pol)
%   text direction      left to right
%   font                Noto Serif
%   fallback font       Noto Serif
%   built from          latex/quickQuestions/fractionsQuickQuestions.tex
%   vocabulary key      latex/quickQuestions/fractionsQuickQuestions/L2/pol/fractionsQuickQuestions_polishKey.tex
%   tier 3 vocabulary   green   RGB 0 128 0
%   English text        black   RGB 0 0 0
%   Polish text         purple  RGB 112 48 160
%   layout macros       version 3
%
% Please don't edit any information in this comment block - the server
% compares it against the current settings and rebuilds this file when
% they differ, so changes here would be overwritten. However, feel free
% to make updates to the LaTeX below if you want to edit it.
% ================================================================

\documentclass[aspectratio=169]{beamer}
% ================================================================
% Copied from the English sheet this was translated from, so that
% anything it sets up for itself still works here
% ================================================================
\usepackage{amsmath}
\usepackage{amssymb}
\usepackage{tikz}
\usetikzlibrary{decorations.pathreplacing}

% ================================================================
% Quick Questions deck: fractions.
%
% Every question gets two slides: the question on its own for the
% mini whiteboards, then the same question again with the solution
% underneath in red. Within each topic the ten questions get
% gradually harder.
%
% Where a picture helps, the solution carries a bar model, a number
% line or a hundred square drawn in the same red as the working.
% ================================================================

\setbeamertemplate{navigation symbols}{}
\setbeamertemplate{footline}{}

% Topic divider.
\newcommand{\qtopic}[1]{%
  \section{#1}%
  \begin{frame}[plain]
    \vfill
    \begin{center}
      {\usebeamercolor[fg]{structure}\Huge\bfseries #1}
    \end{center}
    \vfill
  \end{frame}%
}

% \qq[<question size>][<solution size>]{<question>}{<solution>}
% The two optional sizes drop the type for the wordier questions and
% the longer solutions, so nothing runs off the slide.
\NewDocumentCommand{\qq}{O{\Huge}O{\LARGE}+m+m}{%
  \begin{frame}
    \vfill
    \begin{center}
      \begin{minipage}{0.92\textwidth}
        \centering #1 #3
      \end{minipage}
    \end{center}
    \vfill
  \end{frame}%
  \begin{frame}
    \vfill
    \begin{center}
      \begin{minipage}{0.92\textwidth}
        \centering #1 #3
        \par\vspace{0.9em}
        {\color{red}#2 #4}
      \end{minipage}
    \end{center}
    \vfill
  \end{frame}%
}

% \wholes[<width of one whole, cm>]{<wholes drawn>}{<parts per whole>}{<parts shaded>}
% A row of unit bars, each cut into the same number of parts, with the
% first few parts shaded. Used for mixed numbers and improper fractions.
\NewDocumentCommand{\wholes}{O{1.8}mmm}{%
  \begin{tikzpicture}[x=1cm,y=1cm,baseline=(current bounding box.center)]
    \pgfmathsetmacro{\bw}{#1}
    \pgfmathsetmacro{\cw}{\bw/#3}
    \foreach \w in {1,...,#2}{%
      \pgfmathsetmacro{\bx}{(\w-1)*(\bw+0.22)}%
      \foreach \i in {1,...,#3}{%
        \pgfmathtruncatemacro{\idx}{(\w-1)*#3+\i}%
        \ifnum\idx>#4\relax
          \draw[red!55!black,thin] ({\bx+(\i-1)*\cw},0) rectangle ({\bx+\i*\cw},0.5);
        \else
          \filldraw[fill=red!25,draw=red!55!black,thin]
            ({\bx+(\i-1)*\cw},0) rectangle ({\bx+\i*\cw},0.5);
        \fi}%
      \draw[red!55!black,thick] (\bx,0) rectangle ({\bx+\bw},0.5);}
  \end{tikzpicture}%
}

% A single bar: \fbar[<width, cm>]{<parts>}{<parts shaded>}
\NewDocumentCommand{\fbar}{O{6}mm}{\wholes[#1]{1}{#2}{#3}}

% Bars stacked one above another with a label down the left, written as
%   \barstack{$\frac12$ & \fbar{2}{1} \\[0.95em] $\frac14$ & \fbar{4}{1}}
% Rows carry an explicit skip because \arraystretch only scales the strut,
% which is not enough to keep stacked \dfrac labels apart.
\newcommand{\barstack}[1]{%
  {\renewcommand{\arraystretch}{1.2}%
   \begin{tabular}{@{}r@{\;\;}l@{}}#1\end{tabular}}%
}

% Stepped working, written in the form
%   3/4 of 40 is 30
% Rows are separated with \\ and the two columns with &.
\newcommand{\steps}[1]{%
  {\renewcommand{\arraystretch}{1.5}%
   \begin{tabular}{@{}r@{\ is\ }l@{}}#1\end{tabular}}%
}

% Chained working, one equation per row, for the algebraic methods.
\newcommand{\work}[1]{%
  \begin{tabular}{@{}r@{\;=\;}l@{}}#1\end{tabular}%
}

% \barmodel{<cells>}{<cells spanned by the lower brace>}{<lower label>}{<upper label>}
% The whole bar is 9cm wide and represents the upper label; the lower
% brace picks out the part of it named by the lower label. Kept
% identical to the percentages deck so the two decks match.
\newcommand{\barmodel}[4]{%
  \begin{tikzpicture}[x=1cm,y=1cm]
    \pgfmathsetmacro{\cw}{9/#1}
    \fill[red!12] (0,0) rectangle ({#2*\cw},0.7);
    \foreach \i in {1,...,#1}
      {\draw[red!60!black] ({(\i-1)*\cw},0) rectangle ({\i*\cw},0.7);}
    \draw[red!60!black,thick] (0,0) rectangle (9,0.7);
    \draw[red,decorate,decoration={brace,amplitude=4pt}] (0,0.9) -- (9,0.9)
      node[midway,above=4pt,red,font=\normalsize]{#4};
    \draw[red,decorate,decoration={brace,mirror,amplitude=4pt}] (0,-0.2) -- ({#2*\cw},-0.2)
      node[midway,below=4pt,red,font=\normalsize]{#3};
  \end{tikzpicture}%
}

% \nline{<divisions>}{<label at 0>}{<label at the far end>}{<marks>}
% draws a number line 10cm long. Marks are placed with \mk (label above)
% and \mkb (label below), each taking how far along the line to go as a
% fraction, so that two close values can sit on opposite sides.
\newcommand{\nline}[4]{%
  \begin{tikzpicture}[x=1cm,y=1cm]
    \draw[red!55!black,thick] (0,0) -- (10,0);
    \foreach \i in {0,...,#1}{%
      \pgfmathsetmacro{\xx}{10*\i/#1}%
      \draw[red!55!black] (\xx,-0.12) -- (\xx,0.12);}
    \node[red!55!black,below=10pt,font=\small] at (0,0) {#2};
    \node[red!55!black,below=10pt,font=\small] at (10,0) {#3};
    #4
  \end{tikzpicture}%
}
\newcommand{\mk}[2]{%
  \fill[red] ({10*(#1)},0) circle (2.4pt);
  \node[red,above=5pt,font=\small] at ({10*(#1)},0) {#2};}
\newcommand{\mkb}[2]{%
  \fill[red] ({10*(#1)},0) circle (2.4pt);
  \node[red,below=5pt,font=\small] at ({10*(#1)},0) {#2};}

% \hgrid{<hundredths shaded>} draws a hundred square, for the
% fraction / decimal / percentage conversions. The grid lines are kept
% much lighter than the shading so the shaded part reads at a glance.
\newcommand{\hgrid}[1]{%
  \begin{tikzpicture}[x=0.33cm,y=0.33cm,baseline=(current bounding box.center)]
    \foreach \i in {1,...,#1}{%
      \pgfmathtruncatemacro{\rr}{floor((\i-1)/10)}%
      \pgfmathtruncatemacro{\cc}{mod(\i-1,10)}%
      \fill[red!35] (\cc,-\rr) rectangle ({\cc+1},{-\rr-1});}
    \draw[red!30,very thin,step=1] (0,0) grid (10,-10);
    \draw[red!60!black,thick] (0,0) rectangle (10,-10);
  \end{tikzpicture}%
}

% Contents-slide bullets, colour-coded by difficulty band:
% green for the first four topics, orange for the next four,
% red for the last three.
\setbeamertemplate{section in toc}{%
  \ifnum\inserttocsectionnumber<5
    \textcolor{green!55!black}{$\bullet$}%
  \else\ifnum\inserttocsectionnumber<9
    \textcolor{orange}{$\bullet$}%
  \else
    \textcolor{red}{$\bullet$}%
  \fi\fi
  ~\inserttocsection%
}

\title{Fractions: Quick Questions}
\subtitle{Mini whiteboard questions, with solutions}
\author{}
\date{}

% Characters this language's font has not got are borrowed from another,
% rather than being silently dropped - see docs/translated-sheet-latex.md
\directlua{%
  if luaotfload and luaotfload.add_fallback then
    luaotfload.add_fallback("eallatinfallback",
      { "Noto Serif:mode=harf;" })
  else
    tex.error("this luaotfload is too old to register a fallback font")
  end
}
\usepackage[english]{babel}
\babelprovide[import]{polish}
\babelfont[polish]{rm}[Renderer=Harfbuzz, BoldFont={Noto Serif}, ItalicFont={Noto Serif}, BoldItalicFont={Noto Serif}, RawFeature={fallback=eallatinfallback}]{Noto Serif}
\babelfont[polish]{sf}[Renderer=Harfbuzz, BoldFont={Noto Serif}, ItalicFont={Noto Serif}, BoldItalicFont={Noto Serif}, RawFeature={fallback=eallatinfallback}]{Noto Serif}
\newcommand{\ealtext}[1]{\foreignlanguage{polish}{#1}}
\newcommand{\ealtextblock}[1]{{\raggedright\foreignlanguage{polish}{#1}\par}}
% ================================================================
% EAL layout helpers
% ================================================================
\usepackage{xcolor}

\definecolor{ealkeycolour}{RGB}{0,128,0}
\definecolor{ealenglish}{RGB}{0,0,0}
\definecolor{ealtwocolour}{RGB}{112,48,160}

\newsavebox{\ealboxen}
\newsavebox{\ealboxtr}
\newlength{\ealglosswd}
\newlength{\ealglossraise}
\setlength{\ealglossraise}{1.15em}

% a tier 3 word, in English and in translation alike
\newcommand{\ealkey}[1]{\textcolor{ealkeycolour}{#1}}
\newcommand{\ealkeytr}[1]{\textcolor{ealkeycolour}{\ealtext{#1}}}

% One translated word above one English word. Built from kernel
% primitives, never a tabular, which would break wherever the sheet
% loads array, booktabs or tabularx. The box is as wide as the wider of
% the two, so a long translation pushes its neighbours aside instead of
% overlapping them, and it claims its full height so a table row or a
% TikZ node makes room for it.
\newcommand{\ealgloss}[2]{%
  \sbox{\ealboxen}{\ealkey{#1}}%
  \sbox{\ealboxtr}{\small\textcolor{ealtwocolour}{\ealtext{#2}}}%
  \setlength{\ealglosswd}{\wd\ealboxen}%
  \ifdim\wd\ealboxtr>\ealglosswd\setlength{\ealglosswd}{\wd\ealboxtr}\fi
  \leavevmode
  \makebox[\ealglosswd][c]{%
    \makebox[0pt][c]{\usebox{\ealboxen}}%
    \makebox[0pt][c]{\raisebox{\ealglossraise}{\usebox{\ealboxtr}}}%
  }%
}

% opens up line spacing around a block of glosses, so the raised
% translations sit clear of the line above
\newenvironment{ealglossed}{\par\linespread{2.1}\selectfont\ignorespaces}{\par}

% a whole translated sentence above its English counterpart. block level,
% so it does not belong inside a tabular cell
\newcommand{\ealpara}[2]{%
  \par\addvspace{0.5\baselineskip}%
  {\color{ealtwocolour}\ealtextblock{#1}}%
  \penalty200
  {\color{ealenglish}#2\par}%
  \addvspace{0.5\baselineskip}%
}

\begin{document}

\frame{\titlepage}

\begin{frame}{Contents}
  \small
  \begin{ealglossed}
  \tableofcontents
  \end{ealglossed}
\end{frame}

%================================================================
\section{\ealgloss{Simplifying}{Upraszczanie} \ealgloss{fractions}{ułamki}}
\begin{frame}[plain]
  \vfill
  \begin{center}
    {\usebeamercolor[fg]{structure}\Huge\bfseries
     \begin{ealglossed}
     \ealgloss{Simplifying}{Upraszczanie} \ealgloss{fractions}{ułamki}
     \end{ealglossed}
    }
  \end{center}
  \vfill
\end{frame}
%================================================================

\qq[\Huge][\large]{\begin{ealglossed}\ealgloss{Simplify}{Uprość} $\dfrac{4}{8}$\end{ealglossed}}
  {\barstack{$\dfrac{4}{8}$ & \fbar{8}{4} \\[0.95em]
             $\dfrac{1}{2}$ & \fbar{2}{1}}
   \par\vspace{0.6em}
   $\dfrac{4}{8}=\dfrac{4\div 4}{8\div 4}=\dfrac{1}{2}$}

\qq[\Huge][\large]{\begin{ealglossed}\ealgloss{Simplify}{Uprość} $\dfrac{6}{9}$\end{ealglossed}}
  {\barstack{$\dfrac{6}{9}$ & \fbar{9}{6} \\[0.95em]
             $\dfrac{2}{3}$ & \fbar{3}{2}}
   \par\vspace{0.6em}
   $\dfrac{6}{9}=\dfrac{6\div 3}{9\div 3}=\dfrac{2}{3}$}

\qq[\Huge][\large]{\begin{ealglossed}\ealgloss{Simplify}{Uprość} $\dfrac{12}{20}$\end{ealglossed}}
  {\barstack{$\dfrac{12}{20}$ & \fbar{20}{12} \\[0.95em]
             $\dfrac{3}{5}$ & \fbar{5}{3}}
   \par\vspace{0.6em}
   $\dfrac{12}{20}=\dfrac{12\div 4}{20\div 4}=\dfrac{3}{5}$}

\qq{\begin{ealglossed}\ealgloss{Simplify}{Uprość} $\dfrac{24}{36}$\end{ealglossed}}
  {$\dfrac{24}{36}=\dfrac{24\div 12}{36\div 12}=\dfrac{2}{3}$}

\qq{\begin{ealglossed}\ealgloss{Simplify}{Uprość} $\dfrac{45}{60}$\end{ealglossed}}
  {$\dfrac{45}{60}=\dfrac{45\div 15}{60\div 15}=\dfrac{3}{4}$}

\qq{\begin{ealglossed}\ealgloss{Simplify}{Uprość} $\dfrac{42}{70}$\end{ealglossed}}
  {$\dfrac{42}{70}=\dfrac{42\div 14}{70\div 14}=\dfrac{3}{5}$}

\qq[\Huge][\Large]{\begin{ealglossed}\ealgloss{Simplify}{Uprość} $\dfrac{84}{144}$\end{ealglossed}}
  {$\dfrac{84}{144}=\dfrac{84\div 12}{144\div 12}=\dfrac{7}{12}$}

\qq[\Huge][\Large]{\begin{ealglossed}\ealgloss{Simplify}{Uprość} $\dfrac{91}{117}$\end{ealglossed}}
  {$91=7\times 13$ \qquad $117=9\times 13$\\[0.4em]
   $\dfrac{91}{117}=\dfrac{91\div 13}{117\div 13}=\dfrac{7}{9}$}

\qq[\Huge][\Large]{\begin{ealglossed}\ealgloss{Simplify}{Uprość} $\dfrac{315}{420}$\end{ealglossed}}
  {Both divide by $105$:\\[0.4em]
   $\dfrac{315}{420}=\dfrac{315\div 105}{420\div 105}=\dfrac{3}{4}$}

\qq[\LARGE][\Large]{\begin{ealglossed}
\ealgloss{Simplify}{Uprość} $\dfrac{1001}{1309}$\\[0.4em]
\normalsize (Hint: both numbers are \ealgloss{multiples}{wielokrotności} of $7$ and of $11$.)
\end{ealglossed}}
  {$1001=7\times 11\times 13$ \qquad $1309=7\times 11\times 17$\\[0.4em]
   Divide both by $77$:\\[0.3em]
   $\dfrac{1001}{1309}=\dfrac{13}{17}$}

%================================================================
\section{Ordering \ealgloss{fractions}{ułamki}}
\begin{frame}[plain]
  \vfill
  \begin{center}
    {\usebeamercolor[fg]{structure}\Huge\bfseries
     \begin{ealglossed}
     Ordering \ealgloss{fractions}{ułamki}
     \end{ealglossed}
    }
  \end{center}
  \vfill
\end{frame}
%================================================================

\qq[\LARGE][\large]{Put these in order, smallest first:\\[0.6em]
$\dfrac{1}{4}\qquad \dfrac{1}{2}\qquad \dfrac{1}{8}$}
  {\barstack{$\dfrac18$ & \fbar{8}{1} \\[0.95em]
             $\dfrac14$ & \fbar{4}{1} \\[0.95em]
             $\dfrac12$ & \fbar{2}{1}}
   \par\vspace{0.5em}
   \begin{ealglossed}
   The bigger the \ealgloss{denominator}{mianownik}, the smaller the piece.
   \end{ealglossed}}

\qq[\LARGE][\large]{Put these in order, smallest first:\\[0.6em]
$\dfrac{3}{7}\qquad \dfrac{6}{7}\qquad \dfrac{2}{7}$}
  {\barstack{$\dfrac27$ & \fbar{7}{2} \\[0.95em]
             $\dfrac37$ & \fbar{7}{3} \\[0.95em]
             $\dfrac67$ & \fbar{7}{6}}
   \par\vspace{0.5em}
   \begin{ealglossed}
   Same \ealgloss{denominator}{mianownik}, so just order the \ealgloss{numerators}{liczniki}.
   \end{ealglossed}}

\qq[\LARGE][\large]{Put these in order, smallest first:\\[0.6em]
$\dfrac{1}{2}\qquad \dfrac{3}{4}\qquad \dfrac{5}{8}$}
  {\barstack{$\dfrac12=\dfrac48$ & \fbar{8}{4} \\[0.95em]
             $\dfrac58$ & \fbar{8}{5} \\[0.95em]
             $\dfrac34=\dfrac68$ & \fbar{8}{6}}
   \par\vspace{0.5em}
   $\dfrac12\quad \dfrac58\quad \dfrac34$}

\qq[\LARGE][\large]{Put these in order, smallest first:\\[0.6em]
$\dfrac{2}{3}\qquad \dfrac{5}{6}\qquad \dfrac{3}{4}$}
  {\barstack{$\dfrac23=\dfrac{8}{12}$ & \fbar{12}{8} \\[0.95em]
             $\dfrac34=\dfrac{9}{12}$ & \fbar{12}{9} \\[0.95em]
             $\dfrac56=\dfrac{10}{12}$ & \fbar{12}{10}}
   \par\vspace{0.5em}
   $\dfrac23\quad \dfrac34\quad \dfrac56$}

\qq[\LARGE][\Large]{Put these in order, smallest first:\\[0.6em]
$\dfrac{3}{5}\qquad \dfrac{7}{10}\qquad \dfrac{13}{20}$}
  {\begin{ealglossed}
   \ealgloss{Common denominator}{Wspólny mianownik} $20$:
   \end{ealglossed}\\[0.4em]
   $\dfrac{12}{20}\qquad \dfrac{14}{20}\qquad \dfrac{13}{20}$\\[0.4em]
   $\dfrac{3}{5}\quad \dfrac{13}{20}\quad \dfrac{7}{10}$}

\qq[\LARGE][\Large]{Put these in order, smallest first:\\[0.6em]
$\dfrac{5}{9}\qquad \dfrac{2}{3}\qquad \dfrac{7}{12}$}
  {\begin{ealglossed}
   \ealgloss{Common denominator}{Wspólny mianownik} $36$:
   \end{ealglossed}\\[0.4em]
   $\dfrac{20}{36}\qquad \dfrac{24}{36}\qquad \dfrac{21}{36}$\\[0.4em]
   $\dfrac{5}{9}\quad \dfrac{7}{12}\quad \dfrac{2}{3}$}

\qq[\LARGE][\Large]{Put these in order, \textbf{largest} first:\\[0.6em]
$\dfrac{4}{5}\qquad \dfrac{7}{9}\qquad \dfrac{11}{15}$}
  {\begin{ealglossed}
   \ealgloss{Common denominator}{Wspólny mianownik} $45$:
   \end{ealglossed}\\[0.4em]
   $\dfrac{36}{45}\qquad \dfrac{35}{45}\qquad \dfrac{33}{45}$\\[0.4em]
   $\dfrac{4}{5}\quad \dfrac{7}{9}\quad \dfrac{11}{15}$}

\qq[\LARGE][\Large]{Put these in order, smallest first:\\[0.6em]
$\dfrac{5}{8}\qquad \dfrac{7}{11}\qquad \dfrac{9}{14}$}
  {\begin{ealglossed}
   The \ealgloss{denominators}{mianowniki} share no \ealgloss{factors}{dzielniki}, so divide instead:
   \end{ealglossed}\\[0.4em]
   \work{$\dfrac58$ & $0.625$ \\[0.2em]
         $\dfrac{7}{11}$ & $0.636\ldots$ \\[0.2em]
         $\dfrac{9}{14}$ & $0.642\ldots$}\\[0.4em]
   $\dfrac58\quad \dfrac{7}{11}\quad \dfrac{9}{14}$}

\qq[\LARGE][\large]{Which is closer to $1$: $\dfrac{8}{9}$ or $\dfrac{11}{12}$?\\[0.4em]
\normalsize How do you know?}
  {\barstack{$\dfrac89$ & \fbar{9}{8} \\[0.95em]
             $\dfrac{11}{12}$ & \fbar{12}{11}}
   \par\vspace{0.5em}
   The gaps left are $\dfrac19$ and $\dfrac{1}{12}$.\\[0.3em]
   $\dfrac{1}{12}$ is smaller, so $\dfrac{11}{12}$ is closer to $1$.}

\qq[\LARGE][\Large]{\begin{ealglossed}
Write a \ealgloss{fraction}{ułamek} that lies between\\[0.6em]
$\dfrac{5}{8}$ \quad and \quad $\dfrac{2}{3}$
\end{ealglossed}}
  {In $24$ths they are next-door neighbours:
   $\dfrac{15}{24}$ and $\dfrac{16}{24}$.\\[0.4em]
   \begin{ealglossed}
   Double both \ealgloss{denominators}{mianowniki} to make room:
   \end{ealglossed}\\[0.3em]
   $\dfrac{30}{48}\ <\ \dfrac{31}{48}\ <\ \dfrac{32}{48}$\\[0.4em]
   So $\dfrac{31}{48}$ works.}

%================================================================
\section{\ealgloss{Mixed numbers}{Liczby mieszane} and \ealgloss{improper fractions}{ułamki niewłaściwe}}
\begin{frame}[plain]
  \vfill
  \begin{center}
    {\usebeamercolor[fg]{structure}\Huge\bfseries
     \begin{ealglossed}
     \ealgloss{Mixed numbers}{Liczby mieszane} and \ealgloss{improper fractions}{ułamki niewłaściwe}
     \end{ealglossed}
    }
  \end{center}
  \vfill
\end{frame}
%================================================================

\qq[\Huge][\large]{\begin{ealglossed}
Write $1\dfrac{1}{2}$ as an \ealgloss{improper fraction}{ułamek niewłaściwy}
\end{ealglossed}}
  {\wholes[2]{2}{2}{3}
   \par\vspace{0.6em}
   Three halves are shaded:\\[0.3em]
   $1\dfrac12=\dfrac{2}{2}+\dfrac{1}{2}=\dfrac{3}{2}$}

\qq[\Huge][\large]{\begin{ealglossed}
Write $2\dfrac{3}{4}$ as an \ealgloss{improper fraction}{ułamek niewłaściwy}
\end{ealglossed}}
  {\wholes[2]{3}{4}{11}
   \par\vspace{0.6em}
   $2\dfrac34=\dfrac{8}{4}+\dfrac{3}{4}=\dfrac{11}{4}$}

\qq[\Huge][\large]{\begin{ealglossed}
Write $\dfrac{7}{2}$ as a \ealgloss{mixed number}{liczba mieszana}
\end{ealglossed}}
  {\wholes[2]{4}{2}{7}
   \par\vspace{0.6em}
   \begin{ealglossed}
   $7\div 2=3$ \ealgloss{remainder}{reszta} $1$, so $\dfrac72=3\dfrac12$
   \end{ealglossed}}

\qq[\Huge][\Large]{\begin{ealglossed}
Write $4\dfrac{2}{5}$ as an \ealgloss{improper fraction}{ułamek niewłaściwy}
\end{ealglossed}}
  {$4\times 5=20$, then add the $2$:\\[0.4em]
   $4\dfrac25=\dfrac{22}{5}$}

\qq[\Huge][\large]{\begin{ealglossed}
Write $\dfrac{23}{4}$ as a \ealgloss{mixed number}{liczba mieszana}
\end{ealglossed}}
  {\wholes[1.7]{6}{4}{23}
   \par\vspace{0.6em}
   \begin{ealglossed}
   $23\div 4=5$ \ealgloss{remainder}{reszta} $3$, so $\dfrac{23}{4}=5\dfrac34$
   \end{ealglossed}}

\qq[\Huge][\Large]{\begin{ealglossed}
Write $5\dfrac{7}{8}$ as an \ealgloss{improper fraction}{ułamek niewłaściwy}
\end{ealglossed}}
  {$5\times 8=40$, then add the $7$:\\[0.4em]
   $5\dfrac78=\dfrac{47}{8}$}

\qq[\Huge][\Large]{\begin{ealglossed}
Write $\dfrac{100}{7}$ as a \ealgloss{mixed number}{liczba mieszana}
\end{ealglossed}}
  {$7\times 14=98$, leaving $2$:\\[0.4em]
   $\dfrac{100}{7}=14\dfrac27$}

\qq[\Huge][\Large]{\begin{ealglossed}
Write $12\dfrac{5}{9}$ as an \ealgloss{improper fraction}{ułamek niewłaściwy}
\end{ealglossed}}
  {$12\times 9=108$, then add the $5$:\\[0.4em]
   $12\dfrac59=\dfrac{113}{9}$}

\qq[\Huge][\Large]{\begin{ealglossed}
Write $-3\dfrac{2}{5}$ as an \ealgloss{improper fraction}{ułamek niewłaściwy}
\end{ealglossed}}
  {\begin{ealglossed}
   \ealgloss{Convert}{Zamień} $3\dfrac25$ first, then put the sign back:\\[0.4em]
   $3\times 5+2=17$, so $-3\dfrac25=-\dfrac{17}{5}$
   \end{ealglossed}}

\qq[\LARGE][\Large]{\begin{ealglossed}
The \ealgloss{mixed number}{liczba mieszana} $6\dfrac{a}{7}$ is the same as $\dfrac{45}{7}$.\\[0.6em]
Find $a$.
\end{ealglossed}}
  {$6\dfrac{a}{7}=\dfrac{42+a}{7}$\\[0.4em]
   $42+a=45$, so $a=3$}

%================================================================
\section{Adding and subtracting \ealgloss{fractions}{ułamki}}
\begin{frame}[plain]
  \vfill
  \begin{center}
    {\usebeamercolor[fg]{structure}\Huge\bfseries
     \begin{ealglossed}
     Adding and subtracting \ealgloss{fractions}{ułamki}
     \end{ealglossed}
    }
  \end{center}
  \vfill
\end{frame}
%================================================================

\qq[\Huge][\large]{$\dfrac{1}{2}+\dfrac{1}{4}$}
  {\barstack{$\dfrac12=\dfrac24$ & \fbar{4}{2} \\[0.95em]
             $\dfrac14$ & \fbar{4}{1} \\[0.95em]
             total & \fbar{4}{3}}
   \par\vspace{0.5em}
   $\dfrac24+\dfrac14=\dfrac34$}

\qq[\Huge][\large]{$\dfrac{2}{3}+\dfrac{1}{6}$}
  {\barstack{$\dfrac23=\dfrac46$ & \fbar{6}{4} \\[0.95em]
             $\dfrac16$ & \fbar{6}{1} \\[0.95em]
             total & \fbar{6}{5}}
   \par\vspace{0.5em}
   $\dfrac46+\dfrac16=\dfrac56$}

\qq[\Huge][\large]{$\dfrac{3}{4}-\dfrac{1}{8}$}
  {\barstack{$\dfrac34=\dfrac68$ & \fbar{8}{6} \\[0.95em]
             take off $\dfrac18$ & \fbar{8}{1} \\[0.95em]
             leaves & \fbar{8}{5}}
   \par\vspace{0.5em}
   $\dfrac68-\dfrac18=\dfrac58$}

\qq[\Huge][\large]{$\dfrac{1}{3}+\dfrac{1}{4}$}
  {\barstack{$\dfrac13=\dfrac{4}{12}$ & \fbar{12}{4} \\[0.95em]
             $\dfrac14=\dfrac{3}{12}$ & \fbar{12}{3} \\[0.95em]
             total & \fbar{12}{7}}
   \par\vspace{0.5em}
   $\dfrac{4}{12}+\dfrac{3}{12}=\dfrac{7}{12}$}

\qq[\Huge][\Large]{$\dfrac{3}{5}-\dfrac{1}{4}$}
  {\begin{ealglossed}
   \ealgloss{Common denominator}{Wspólny mianownik} $20$:
   \end{ealglossed}\\[0.4em]
   $\dfrac{12}{20}-\dfrac{5}{20}=\dfrac{7}{20}$}

\qq[\Huge][\Large]{$\dfrac{5}{6}+\dfrac{3}{8}$}
  {\begin{ealglossed}
   \ealgloss{Common denominator}{Wspólny mianownik} $24$:
   \end{ealglossed}\\[0.4em]
   $\dfrac{20}{24}+\dfrac{9}{24}=\dfrac{29}{24}=1\dfrac{5}{24}$}

\qq[\Huge][\Large]{$\dfrac{7}{10}-\dfrac{2}{15}$}
  {\begin{ealglossed}
   \ealgloss{Common denominator}{Wspólny mianownik} $30$:
   \end{ealglossed}\\[0.4em]
   $\dfrac{21}{30}-\dfrac{4}{30}=\dfrac{17}{30}$}

\qq[\Huge][\Large]{$\dfrac{2}{3}+\dfrac{3}{4}+\dfrac{1}{6}$}
  {\begin{ealglossed}
   \ealgloss{Common denominator}{Wspólny mianownik} $12$:
   \end{ealglossed}\\[0.4em]
   $\dfrac{8}{12}+\dfrac{9}{12}+\dfrac{2}{12}=\dfrac{19}{12}=1\dfrac{7}{12}$}

\qq[\Huge][\Large]{$\dfrac{5}{12}-\dfrac{7}{18}$}
  {\begin{ealglossed}
   \ealgloss{Common denominator}{Wspólny mianownik} $36$:
   \end{ealglossed}\\[0.4em]
   $\dfrac{15}{36}-\dfrac{14}{36}=\dfrac{1}{36}$}

\qq[\LARGE][\Large]{What must be added to $\dfrac{5}{9}$\\[0.4em]
to make $\dfrac{7}{8}$?}
  {The missing amount is $\dfrac78-\dfrac59$.\\[0.4em]
   \begin{ealglossed}
   \ealgloss{Common denominator}{Wspólny mianownik} $72$:
   \end{ealglossed}\\[0.3em]
   $\dfrac{63}{72}-\dfrac{40}{72}=\dfrac{23}{72}$}

%================================================================
\section{\ealgloss{Converting}{Zamienianie} \ealgloss{fractions}{ułamki}, \ealgloss{decimals}{ułamki dziesiętne} and \ealgloss{percentages}{procenty}}
\begin{frame}[plain]
  \vfill
  \begin{center}
    {\usebeamercolor[fg]{structure}\Huge\bfseries
     \begin{ealglossed}
     \ealgloss{Converting}{Zamienianie} \ealgloss{fractions}{ułamki}, \ealgloss{decimals}{ułamki dziesiętne} and \ealgloss{percentages}{procenty}
     \end{ealglossed}
    }
  \end{center}
  \vfill
\end{frame}
%================================================================

\qq[\LARGE][\large]{\begin{ealglossed}
Write $\dfrac{1}{2}$ as a \ealgloss{decimal}{ułamek dziesiętny} and as a \ealgloss{percentage}{procent}
\end{ealglossed}}
  {\hgrid{50}
   \par\vspace{0.5em}
   Half the square is $50$ out of $100$:\\[0.3em]
   $\dfrac12=\dfrac{50}{100}=0.5=50\%$}

\qq[\LARGE][\large]{\begin{ealglossed}
Write $0.25$ as a \ealgloss{fraction}{ułamek} and as a \ealgloss{percentage}{procent}
\end{ealglossed}}
  {\hgrid{25}
   \par\vspace{0.5em}
   $0.25=\dfrac{25}{100}=\dfrac14=25\%$}

\qq[\LARGE][\large]{\begin{ealglossed}
Write $30\%$ as a \ealgloss{decimal}{ułamek dziesiętny} and as a \ealgloss{fraction}{ułamek} in its \ealgloss{simplest form}{najprostsza postać}
\end{ealglossed}}
  {\hgrid{30}
   \par\vspace{0.5em}
   $30\%=\dfrac{30}{100}=0.3=\dfrac{3}{10}$}

\qq[\LARGE][\Large]{\begin{ealglossed}
Write $\dfrac{3}{4}$ as a \ealgloss{percentage}{procent}, and $0.6$ as a \ealgloss{percentage}{procent}
\end{ealglossed}}
  {$\dfrac34=\dfrac{75}{100}=75\%$\\[0.4em]
   $0.6=\dfrac{60}{100}=60\%$}

\qq[\LARGE][\large]{\begin{ealglossed}
Write $\dfrac{7}{20}$ as a \ealgloss{decimal}{ułamek dziesiętny} and as a \ealgloss{percentage}{procent}
\end{ealglossed}}
  {\hgrid{35}
   \par\vspace{0.5em}
   Multiply top and bottom by $5$:\\[0.3em]
   $\dfrac{7}{20}=\dfrac{35}{100}=0.35=35\%$}

\qq[\LARGE][\large]{\begin{ealglossed}
Write $0.08$ as a \ealgloss{percentage}{procent} and as a \ealgloss{fraction}{ułamek} in its \ealgloss{simplest form}{najprostsza postać}
\end{ealglossed}}
  {\hgrid{8}
   \par\vspace{0.5em}
   $0.08=\dfrac{8}{100}=8\%=\dfrac{2}{25}$}

\qq[\LARGE][\Large]{\begin{ealglossed}
Write $\dfrac{3}{8}$ as a \ealgloss{decimal}{ułamek dziesiętny} and as a \ealgloss{percentage}{procent}
\end{ealglossed}}
  {$3\div 8=0.375$\\[0.4em]
   $0.375\times 100=37.5\%$}

\qq[\LARGE][\Large]{\begin{ealglossed}
Write $12.5\%$ as a \ealgloss{fraction}{ułamek} in its \ealgloss{simplest form}{najprostsza postać}
\end{ealglossed}}
  {$12.5\%=\dfrac{12.5}{100}=\dfrac{125}{1000}$\\[0.4em]
   Divide both by $125$: \ $\dfrac{1}{8}$}

\qq[\LARGE][\Large]{\begin{ealglossed}
Write $\dfrac{5}{6}$ as a \ealgloss{decimal}{ułamek dziesiętny}, and as a \ealgloss{percentage}{procent} to $1$ \ealgloss{decimal place}{miejsce dziesiętne}
\end{ealglossed}}
  {$5\div 6=0.8333\ldots$\\[0.4em]
   $0.8333\ldots\times 100=83.3\%$ (1 d.p.)}

\qq[\LARGE][\Large]{\begin{ealglossed}
Write $137.5\%$ as a \ealgloss{mixed number}{liczba mieszana} in its \ealgloss{simplest form}{najprostsza postać}
\end{ealglossed}}
  {$137.5\%=1.375=1\dfrac{375}{1000}$\\[0.4em]
   Divide both by $125$: \ $1\dfrac{3}{8}$}

%================================================================
\section{Ordering \ealgloss{fractions}{ułamki}, \ealgloss{decimals}{ułamki dziesiętne} and \ealgloss{percentages}{procenty}}
\begin{frame}[plain]
  \vfill
  \begin{center}
    {\usebeamercolor[fg]{structure}\Huge\bfseries
     \begin{ealglossed}
     Ordering \ealgloss{fractions}{ułamki}, \ealgloss{decimals}{ułamki dziesiętne} and \ealgloss{percentages}{procenty}
     \end{ealglossed}
    }
  \end{center}
  \vfill
\end{frame}
%================================================================

\qq[\LARGE][\large]{Put these in order, smallest first:\\[0.6em]
$0.5\qquad \dfrac{1}{4}\qquad 60\%$}
  {\nline{10}{$0$}{$1$}{\mk{0.25}{$\frac14$}\mkb{0.5}{$0.5$}\mk{0.6}{$60\%$}}
   \par\vspace{0.6em}
   $\dfrac14=0.25$ \qquad $60\%=0.6$\\[0.3em]
   $\dfrac14\quad 0.5\quad 60\%$}

\qq[\LARGE][\large]{Put these in order, smallest first:\\[0.6em]
$0.7\qquad \dfrac{3}{4}\qquad 70\%$}
  {\nline{10}{$0$}{$1$}{\mkb{0.7}{$0.7=70\%$}\mk{0.75}{$\frac34$}}
   \par\vspace{0.6em}
   $0.7$ and $70\%$ are equal, and $\dfrac34=0.75$.\\[0.3em]
   $0.7=70\%$, \ then \ $\dfrac34$}

\qq[\LARGE][\Large]{Put these in order, smallest first:\\[0.6em]
$\dfrac{2}{5}\qquad 45\%\qquad 0.35$}
  {\begin{ealglossed}
   As \ealgloss{decimals}{ułamki dziesiętne}: \ $0.4$ \quad $0.45$ \quad $0.35$
   \end{ealglossed}\\[0.4em]
   $0.35\quad \dfrac25\quad 45\%$}

\qq[\LARGE][\Large]{Put these in order, smallest first:\\[0.6em]
$0.62\qquad \dfrac{5}{8}\qquad 63\%$}
  {\begin{ealglossed}
   As \ealgloss{decimals}{ułamki dziesiętne}: \ $0.62$ \quad $0.625$ \quad $0.63$
   \end{ealglossed}\\[0.4em]
   $0.62\quad \dfrac58\quad 63\%$}

\qq[\LARGE][\Large]{Put these in order, \textbf{largest} first:\\[0.6em]
$\dfrac{7}{8}\qquad 0.87\qquad 88\%$}
  {\begin{ealglossed}
   As \ealgloss{decimals}{ułamki dziesiętne}: \ $0.875$ \quad $0.87$ \quad $0.88$
   \end{ealglossed}\\[0.4em]
   $88\%\quad \dfrac78\quad 0.87$}

\qq[\LARGE][\Large]{Put these in order, smallest first:\\[0.6em]
$\dfrac{1}{3}\qquad 0.3\qquad 33\%$}
  {\begin{ealglossed}
   As \ealgloss{decimals}{ułamki dziesiętne}: \ $0.333\ldots$ \quad $0.3$ \quad $0.33$
   \end{ealglossed}\\[0.4em]
   $0.3\quad 33\%\quad \dfrac13$}

\qq[\LARGE][\Large]{Put these in order, smallest first:\\[0.6em]
$0.045\qquad 4.5\%\qquad \dfrac{1}{20}$}
  {\begin{ealglossed}
   As \ealgloss{decimals}{ułamki dziesiętne}: \ $0.045$ \quad $0.045$ \quad $0.05$
   \end{ealglossed}\\[0.4em]
   $0.045=4.5\%$, \ then \ $\dfrac{1}{20}$}

\qq[\LARGE][\Large]{Put these in order, smallest first:\\[0.6em]
$\dfrac{5}{6}\qquad 83\%\qquad 0.835$}
  {\begin{ealglossed}
   As \ealgloss{decimals}{ułamki dziesiętne}: \ $0.8333\ldots$ \quad $0.83$ \quad $0.835$
   \end{ealglossed}\\[0.4em]
   $83\%\quad \dfrac56\quad 0.835$}

\qq[\LARGE][\Large]{Put these in order, smallest first:\\[0.6em]
$\dfrac{2}{7}\qquad 0.28\qquad 29\%$}
  {\begin{ealglossed}
   As \ealgloss{decimals}{ułamki dziesiętne}: \ $0.2857\ldots$ \quad $0.28$ \quad $0.29$
   \end{ealglossed}\\[0.4em]
   $0.28\quad \dfrac27\quad 29\%$}

\qq[\LARGE][\Large]{Put these in order, smallest first:\\[0.6em]
$\dfrac{3}{8}\qquad 0.38\qquad 37.5\%\qquad \dfrac{3}{7}$\\[0.5em]
\normalsize What do you notice?}
  {\begin{ealglossed}
   As \ealgloss{decimals}{ułamki dziesiętne}: \ $0.375$ \quad $0.38$ \quad $0.375$ \quad $0.4285\ldots$
   \end{ealglossed}\\[0.4em]
   $\dfrac38=37.5\%$, \ then \ $0.38$, \ then \ $\dfrac37$\\[0.4em]
   $\dfrac38$ and $37.5\%$ are the same number written two ways.}

%================================================================
\section{\ealgloss{Fractions}{Ułamki} of amounts (no calculator)}
\begin{frame}[plain]
  \vfill
  \begin{center}
    {\usebeamercolor[fg]{structure}\Huge\bfseries
     \begin{ealglossed}
     \ealgloss{Fractions}{Ułamki} of amounts (no calculator)
     \end{ealglossed}
    }
  \end{center}
  \vfill
\end{frame}
%================================================================

\qq[\Huge][\Large]{$\dfrac{1}{2}$ of $40$}
  {\barmodel{2}{1}{$\frac12$ of $40$ is $20$}{$40$}
   \par\vspace{0.7em}
   $40\div 2=20$}

\qq[\Huge][\Large]{$\dfrac{1}{4}$ of $60$}
  {\barmodel{4}{1}{$\frac14$ of $60$ is $15$}{$60$}
   \par\vspace{0.7em}
   $60\div 4=15$}

\qq[\Huge][\Large]{$\dfrac{1}{5}$ of $45$}
  {\barmodel{5}{1}{$\frac15$ of $45$ is $9$}{$45$}
   \par\vspace{0.7em}
   $45\div 5=9$}

\qq[\Huge][\Large]{$\dfrac{3}{4}$ of $40$}
  {\barmodel{4}{3}{$\frac34$ of $40$ is $30$}{$40$}
   \par\vspace{0.7em}
   \steps{$\frac14$ of $40$ & $10$ \\ $\frac34$ of $40$ & $30$}}

\qq[\Huge][\Large]{$\dfrac{2}{3}$ of $36$}
  {\barmodel{3}{2}{$\frac23$ of $36$ is $24$}{$36$}
   \par\vspace{0.7em}
   \steps{$\frac13$ of $36$ & $12$ \\ $\frac23$ of $36$ & $24$}}

\qq[\Huge][\Large]{$\dfrac{3}{5}$ of $45$}
  {\barmodel{5}{3}{$\frac35$ of $45$ is $27$}{$45$}
   \par\vspace{0.7em}
   \steps{$\frac15$ of $45$ & $9$ \\ $\frac35$ of $45$ & $27$}}

\qq[\Huge][\Large]{$\dfrac{5}{8}$ of $72$}
  {\barmodel{8}{5}{$\frac58$ of $72$ is $45$}{$72$}
   \par\vspace{0.7em}
   \steps{$\frac18$ of $72$ & $9$ \\ $\frac58$ of $72$ & $45$}}

\qq[\Huge][\Large]{$\dfrac{2}{7}$ of $350$}
  {\barmodel{7}{2}{$\frac27$ of $350$ is $100$}{$350$}
   \par\vspace{0.7em}
   \steps{$\frac17$ of $350$ & $50$ \\ $\frac27$ of $350$ & $100$}}

\qq[\LARGE][\large]{$\dfrac{5}{6}$ of a number is $45$.\\[0.4em]
What is the number?}
  {\barmodel{6}{5}{$\frac56$ of it is $45$}{the number}
   \par\vspace{0.7em}
   \steps{$\frac56$ of it & $45$ \\ $\frac16$ of it & $9$ \\ the number & $54$}}

\qq[\LARGE][\Large]{$\dfrac{3}{8}$ of $\dfrac{2}{5}$ of $200$}
  {Work from the inside out:\\[0.4em]
   \steps{$\frac25$ of $200$ & $80$ \\ $\frac18$ of $80$ & $10$ \\ $\frac38$ of $80$ & $30$}}

%================================================================
\section{\ealgloss{Fractions}{Ułamki} of amounts (calculator)}
\begin{frame}[plain]
  \vfill
  \begin{center}
    {\usebeamercolor[fg]{structure}\Huge\bfseries
     \begin{ealglossed}
     \ealgloss{Fractions}{Ułamki} of amounts (calculator)
     \end{ealglossed}
    }
  \end{center}
  \vfill
\end{frame}
%================================================================

\qq[\Huge][\Large]{$\dfrac{1}{4}$ of $386$}
  {$386\div 4=96.5$}

\qq[\Huge][\Large]{$\dfrac{3}{5}$ of $245$}
  {\steps{$\frac15$ of $245$ & $49$ \\ $\frac35$ of $245$ & $147$}}

\qq[\Huge][\Large]{$\dfrac{2}{7}$ of $483$}
  {\steps{$\frac17$ of $483$ & $69$ \\ $\frac27$ of $483$ & $138$}}

\qq[\Huge][\Large]{$\dfrac{3}{16}$ of $1024$}
  {\steps{$\frac{1}{16}$ of $1024$ & $64$ \\ $\frac{3}{16}$ of $1024$ & $192$}}

\qq[\Huge][\Large]{$\dfrac{5}{9}$ of $738$}
  {\steps{$\frac19$ of $738$ & $82$ \\ $\frac59$ of $738$ & $410$}}

\qq[\LARGE][\Large]{$\dfrac{7}{8}$ of $\pounds 312.40$}
  {\steps{$\frac18$ of it & \pounds $39.05$ \\ $\frac78$ of it & \pounds $273.35$}}

\qq[\Huge][\Large]{$\dfrac{11}{13}$ of $2405$}
  {\steps{$\frac{1}{13}$ of $2405$ & $185$ \\ $\frac{11}{13}$ of $2405$ & $2035$}}

\qq[\LARGE][\Large]{$\dfrac{4}{7}$ of $3.5$ kg.\\[0.4em]
Give your answer in grams.}
  {\steps{$\frac17$ of $3.5$ kg & $0.5$ kg \\ $\frac47$ of $3.5$ kg & $2$ kg}
   \par\vspace{0.4em}
   $2$ kg $=2000$ g}

\qq[\LARGE][\large]{$\dfrac{5}{8}$ of a number is $215$.\\[0.4em]
What is the number?}
  {\barmodel{8}{5}{$\frac58$ of it is $215$}{the number}
   \par\vspace{0.7em}
   \steps{$\frac58$ of it & $215$ \\ $\frac18$ of it & $43$ \\ the number & $344$}}

\qq[\LARGE][\Large]{$\dfrac{7}{12}$ of $4$ hours $48$ minutes.\\[0.4em]
Give your answer in hours and minutes.}
  {$4$ h $48$ min $=288$ minutes\\[0.4em]
   \steps{$\frac{1}{12}$ of $288$ & $24$ \\ $\frac{7}{12}$ of $288$ & $168$}
   \par\vspace{0.4em}
   $168$ min $=2$ hours $48$ minutes}

%================================================================
\section{Four \ealgloss{operations}{działania} with \ealgloss{mixed numbers}{liczby mieszane} and \ealgloss{negatives}{ujemne}}
\begin{frame}[plain]
  \vfill
  \begin{center}
    {\usebeamercolor[fg]{structure}\Huge\bfseries
     \begin{ealglossed}
     Four \ealgloss{operations}{działania} with \ealgloss{mixed numbers}{liczby mieszane} and \ealgloss{negatives}{ujemne}
     \end{ealglossed}
    }
  \end{center}
  \vfill
\end{frame}
%================================================================

\qq[\Huge][\Large]{$1\dfrac{1}{2}+2\dfrac{1}{4}$}
  {\begin{ealglossed}
   Add the wholes, then the \ealgloss{fractions}{ułamki}:
   \end{ealglossed}\\[0.4em]
   $1+2=3$ \qquad $\dfrac12+\dfrac14=\dfrac34$\\[0.4em]
   $=3\dfrac34$}

\qq[\Huge][\Large]{$3\dfrac{4}{5}-1\dfrac{2}{5}$}
  {$3-1=2$ \qquad $\dfrac45-\dfrac25=\dfrac25$\\[0.4em]
   $=2\dfrac25$}

\qq[\LARGE][\Large]{$2\dfrac{1}{2}+\left(-1\dfrac{1}{4}\right)$}
  {\begin{ealglossed}
   Adding a \ealgloss{negative}{ujemny} is subtracting:
   \end{ealglossed}\\[0.4em]
   $2\dfrac12-1\dfrac14=1\dfrac14$}

\qq[\Huge][\Large]{$4\dfrac{1}{6}-2\dfrac{3}{4}$}
  {\begin{ealglossed}
   \ealgloss{Improper fractions}{Ułamki niewłaściwe}, \ealgloss{common denominator}{wspólny mianownik} $12$:
   \end{ealglossed}\\[0.4em]
   $\dfrac{25}{6}-\dfrac{11}{4}=\dfrac{50}{12}-\dfrac{33}{12}=\dfrac{17}{12}=1\dfrac{5}{12}$}

\qq[\Huge][\Large]{$2\dfrac{1}{2}\times 4$}
  {Four lots of $2\dfrac12$:\\[0.4em]
   $2\dfrac12+2\dfrac12+2\dfrac12+2\dfrac12=10$\\[0.4em]
   or $\dfrac52\times 4=\dfrac{20}{2}=10$}

\qq[\Huge][\Large]{$-3\dfrac{1}{3}+1\dfrac{5}{6}$}
  {$-\dfrac{10}{3}+\dfrac{11}{6}=-\dfrac{20}{6}+\dfrac{11}{6}$\\[0.4em]
   $=-\dfrac{9}{6}=-\dfrac32=-1\dfrac12$}

\qq[\Huge][\Large]{$1\dfrac{1}{2}\times 2\dfrac{2}{3}$}
  {\begin{ealglossed}
   \ealgloss{Improper fractions}{Ułamki niewłaściwe} first:
   \end{ealglossed}\\[0.4em]
   $\dfrac32\times\dfrac83=\dfrac{24}{6}=4$}

\qq[\Huge][\Large]{$3\dfrac{3}{4}\div 1\dfrac{1}{2}$}
  {\begin{ealglossed}
   \ealgloss{Improper fractions}{Ułamki niewłaściwe}, then multiply by the \ealgloss{reciprocal}{odwrotność}:
   \end{ealglossed}\\[0.4em]
   $\dfrac{15}{4}\div\dfrac32=\dfrac{15}{4}\times\dfrac23=\dfrac{30}{12}=\dfrac52=2\dfrac12$}

\qq[\Huge][\Large]{$-5\dfrac{1}{3}\div \dfrac{2}{3}$}
  {$-\dfrac{16}{3}\div\dfrac23=-\dfrac{16}{3}\times\dfrac32=-\dfrac{48}{6}=-8$}

\qq[\LARGE][\Large]{$-2\dfrac{1}{4}\times\left(-1\dfrac{1}{3}\right)+1\dfrac{1}{2}$}
  {\begin{ealglossed}
   Multiply first; two \ealgloss{negatives}{ujemne} give a \ealgloss{positive}{liczba dodatnia}:
   \end{ealglossed}\\[0.4em]
   $-\dfrac94\times-\dfrac43=\dfrac{36}{12}=3$\\[0.4em]
   $3+1\dfrac12=4\dfrac12$}

%================================================================
\section{Word problems with \ealgloss{mixed numbers}{liczby mieszane}}
\begin{frame}[plain]
  \vfill
  \begin{center}
    {\usebeamercolor[fg]{structure}\Huge\bfseries
     \begin{ealglossed}
     Word problems with \ealgloss{mixed numbers}{liczby mieszane}
     \end{ealglossed}
    }
  \end{center}
  \vfill
\end{frame}
%================================================================

\qq[\LARGE][\large]{A recipe needs $1\dfrac{1}{2}$ cups of flour.\\[0.4em]
Alex makes a double batch.\\[0.4em]
How much flour is needed?}
  {\wholes[2]{3}{2}{6}
   \par\vspace{0.6em}
   $1\dfrac12\times 2=\dfrac32\times 2=\dfrac62=3$ cups}

\qq[\LARGE][\large]{A plank is $3\dfrac{1}{4}$\,m long.\\[0.4em]
A piece $1\dfrac{1}{2}$\,m long is cut off.\\[0.4em]
How much plank is left?}
  {\barstack{$3\frac14$\,m & \wholes[1.5]{4}{4}{13} \\[0.4em]
             cut off $1\frac12$\,m & \wholes[1.5]{2}{4}{6}}
   \par\vspace{0.5em}
   $\dfrac{13}{4}-\dfrac{6}{4}=\dfrac{7}{4}=1\dfrac34$\,m left}

\qq[\LARGE][\Large]{Priya has two ribbons, $2\dfrac{1}{3}$\,m and $1\dfrac{5}{6}$\,m long.\\[0.4em]
What is their total length?}
  {\begin{ealglossed}
   \ealgloss{Common denominator}{Wspólny mianownik} $6$:
   \end{ealglossed}\\[0.4em]
   $\dfrac{14}{6}+\dfrac{11}{6}=\dfrac{25}{6}=4\dfrac16$\,m}

\qq[\LARGE][\Large]{Sam runs $2\dfrac{3}{4}$ miles on Monday and $3\dfrac{1}{2}$ miles on Tuesday.\\[0.4em]
How much further did he run on Tuesday?}
  {$3\dfrac12-2\dfrac34=\dfrac{14}{4}-\dfrac{11}{4}=\dfrac34$\\[0.4em]
   $\dfrac34$ of a mile further}

\qq[\LARGE][\large]{A tin holds $4\dfrac{1}{2}$ litres of paint.\\[0.4em]
Five jugs of $\dfrac{3}{4}$ of a litre are poured out.\\[0.4em]
How much paint is left?}
  {\barstack{tin, $4\frac12$\,l & \wholes[1.3]{5}{4}{18} \\[0.4em]
             poured, $3\frac34$\,l & \wholes[1.3]{4}{4}{15}}
   \par\vspace{0.5em}
   $5\times\dfrac34=\dfrac{15}{4}$, \ and \
   $\dfrac{18}{4}-\dfrac{15}{4}=\dfrac34$ of a litre left}

\qq[\LARGE][\Large]{A wall is $7\dfrac{1}{2}$\,m wide.\\[0.4em]
Panels are $1\dfrac{1}{4}$\,m wide.\\[0.4em]
How many panels fit across the wall?}
  {$7\dfrac12\div 1\dfrac14=\dfrac{15}{2}\div\dfrac54$\\[0.4em]
   $=\dfrac{15}{2}\times\dfrac45=\dfrac{60}{10}=6$ panels}

\qq[\LARGE][\large]{A bottle holds $1\dfrac{1}{3}$ litres of water.\\[0.4em]
Jo drinks $\dfrac{3}{8}$ of it.\\[0.4em]
How much water does Jo drink?}
  {\barmodel{8}{3}{Jo drinks $\frac12$ litre}{$1\frac13$ litres}
   \par\vspace{0.7em}
   $\dfrac38\times\dfrac43=\dfrac{12}{24}=\dfrac12$ of a litre}

\qq[\LARGE][\large]{A journey takes $2\dfrac{1}{4}$ hours.\\[0.4em]
Maria has driven $\dfrac{2}{3}$ of the way.\\[0.4em]
How long is left, in minutes?}
  {\barmodel{3}{2}{driven $\frac23$, so $\frac13$ is left}{$2\frac14$ hours $=135$ min}
   \par\vspace{0.7em}
   \steps{$\frac13$ of $135$ min & $45$ min}}

\qq[\LARGE][\Large]{A rope $12\dfrac{1}{2}$\,m long is cut into pieces $1\dfrac{3}{4}$\,m long.\\[0.4em]
How many whole pieces are there, and how much rope is left over?}
  {$\dfrac{25}{2}\div\dfrac74=\dfrac{25}{2}\times\dfrac47=\dfrac{50}{7}=7\dfrac17$\\[0.4em]
   $7$ whole pieces use $7\times 1\dfrac34=12\dfrac14$\,m,\\[0.3em]
   leaving $12\dfrac12-12\dfrac14=\dfrac14$\,m of rope.}

\qq[\large][\Large]{A recipe for $4$ people uses $2\dfrac{1}{2}$ cups of rice and $1\dfrac{3}{4}$ litres of stock.\\[0.5em]
How much rice and stock are needed for $10$ people?}
  {$10$ people is $2\dfrac12$ times as many, so multiply by $\dfrac52$:\\[0.4em]
   Rice: $\dfrac52\times\dfrac52=\dfrac{25}{4}=6\dfrac14$ cups\\[0.3em]
   Stock: $\dfrac74\times\dfrac52=\dfrac{35}{8}=4\dfrac38$ litres}

%================================================================
\section{\ealgloss{Recurring decimals}{Ułamki okresowe} to \ealgloss{fractions}{ułamki}}
\begin{frame}[plain]
  \vfill
  \begin{center}
    {\usebeamercolor[fg]{structure}\Huge\bfseries
     \begin{ealglossed}
     \ealgloss{Recurring decimals}{Ułamki okresowe} to \ealgloss{fractions}{ułamki}
     \end{ealglossed}
    }
  \end{center}
  \vfill
\end{frame}
%================================================================

\qq[\LARGE][\Large]{\begin{ealglossed}
Use an \ealgloss{algebraic}{algebraiczny} method to write\\[0.5em]
$0.\dot{3}$ \quad as a \ealgloss{fraction}{ułamek}
\end{ealglossed}}
  {\work{$x$ & $0.333\ldots$ \\[0.2em]
         $10x$ & $3.333\ldots$ \\[0.2em]
         $9x$ & $3$ \\[0.2em]
         $x$ & $\dfrac39=\dfrac13$}}

\qq[\LARGE][\Large]{\begin{ealglossed}
Use an \ealgloss{algebraic}{algebraiczny} method to write\\[0.5em]
$0.\dot{7}$ \quad as a \ealgloss{fraction}{ułamek}
\end{ealglossed}}
  {\work{$x$ & $0.777\ldots$ \\[0.2em]
         $10x$ & $7.777\ldots$ \\[0.2em]
         $9x$ & $7$ \\[0.2em]
         $x$ & $\dfrac79$}}

\qq[\LARGE][\Large]{\begin{ealglossed}
Use an \ealgloss{algebraic}{algebraiczny} method to write\\[0.5em]
$0.\dot{4}\dot{5}$ \quad as a \ealgloss{fraction}{ułamek}
\end{ealglossed}}
  {Two digits repeat, so multiply by $100$:\\[0.4em]
   \work{$x$ & $0.4545\ldots$ \\[0.2em]
         $100x$ & $45.4545\ldots$ \\[0.2em]
         $99x$ & $45$ \\[0.2em]
         $x$ & $\dfrac{45}{99}=\dfrac{5}{11}$}}

\qq[\LARGE][\Large]{\begin{ealglossed}
Use an \ealgloss{algebraic}{algebraiczny} method to write\\[0.5em]
$0.\dot{1}\dot{2}$ \quad as a \ealgloss{fraction}{ułamek} in its \ealgloss{simplest form}{najprostsza postać}
\end{ealglossed}}
  {\work{$x$ & $0.1212\ldots$ \\[0.2em]
         $100x$ & $12.1212\ldots$ \\[0.2em]
         $99x$ & $12$ \\[0.2em]
         $x$ & $\dfrac{12}{99}=\dfrac{4}{33}$}}

\qq[\LARGE][\Large]{\begin{ealglossed}
Use an \ealgloss{algebraic}{algebraiczny} method to write\\[0.5em]
$0.\dot{2}\dot{7}$ \quad as a \ealgloss{fraction}{ułamek} in its \ealgloss{simplest form}{najprostsza postać}
\end{ealglossed}}
  {\work{$x$ & $0.2727\ldots$ \\[0.2em]
         $100x$ & $27.2727\ldots$ \\[0.2em]
         $99x$ & $27$ \\[0.2em]
         $x$ & $\dfrac{27}{99}=\dfrac{3}{11}$}}

\qq[\LARGE][\Large]{\begin{ealglossed}
Use an \ealgloss{algebraic}{algebraiczny} method to write\\[0.5em]
$0.\dot{1}2\dot{3}$ \quad as a \ealgloss{fraction}{ułamek} in its \ealgloss{simplest form}{najprostsza postać}
\end{ealglossed}}
  {Three digits repeat, so multiply by $1000$:\\[0.4em]
   \work{$x$ & $0.123123\ldots$ \\[0.2em]
         $1000x$ & $123.123123\ldots$ \\[0.2em]
         $999x$ & $123$ \\[0.2em]
         $x$ & $\dfrac{123}{999}=\dfrac{41}{333}$}}

\qq[\LARGE][\Large]{\begin{ealglossed}
Use an \ealgloss{algebraic}{algebraiczny} method to write\\[0.5em]
$0.1\dot{6}$ \quad as a \ealgloss{fraction}{ułamek} in its \ealgloss{simplest form}{najprostsza postać}
\end{ealglossed}}
  {Only the $6$ repeats, so line up $10x$ and $100x$:\\[0.4em]
   \work{$10x$ & $1.666\ldots$ \\[0.2em]
         $100x$ & $16.666\ldots$ \\[0.2em]
         $90x$ & $15$ \\[0.2em]
         $x$ & $\dfrac{15}{90}=\dfrac{1}{6}$}}

\qq[\LARGE][\Large]{\begin{ealglossed}
Use an \ealgloss{algebraic}{algebraiczny} method to write\\[0.5em]
$0.8\dot{3}$ \quad as a \ealgloss{fraction}{ułamek} in its \ealgloss{simplest form}{najprostsza postać}
\end{ealglossed}}
  {\work{$10x$ & $8.333\ldots$ \\[0.2em]
         $100x$ & $83.333\ldots$ \\[0.2em]
         $90x$ & $75$ \\[0.2em]
         $x$ & $\dfrac{75}{90}=\dfrac{5}{6}$}}

\qq[\LARGE][\Large]{\begin{ealglossed}
Use an \ealgloss{algebraic}{algebraiczny} method to write\\[0.5em]
$1.\dot{2}\dot{7}$ \quad as a \ealgloss{mixed number}{liczba mieszana} in its \ealgloss{simplest form}{najprostsza postać}
\end{ealglossed}}
  {\begin{ealglossed}
   Deal with the \ealgloss{decimal}{ułamek dziesiętny} part, then put the $1$ back:
   \end{ealglossed}\\[0.4em]
   \work{$x$ & $0.2727\ldots$ \\[0.2em]
         $99x$ & $27$ \\[0.2em]
         $x$ & $\dfrac{27}{99}=\dfrac{3}{11}$}
   \par\vspace{0.4em}
   $1.\dot{2}\dot{7}=1\dfrac{3}{11}$}

\qq[\LARGE][\Large]{\begin{ealglossed}
Use an \ealgloss{algebraic}{algebraiczny} method to show that\\[0.5em]
$0.\dot{9}=1$
\end{ealglossed}}
  {\work{$x$ & $0.999\ldots$ \\[0.2em]
         $10x$ & $9.999\ldots$ \\[0.2em]
         $9x$ & $9$ \\[0.2em]
         $x$ & $1$}
   \par\vspace{0.4em}
   So $0.\dot{9}$ and $1$ are the same number.}

\end{document}